% prezentace.tex -- Lekce 02: Proměnné a typy
% preamble-prez.tex -- sdílená preamble pro prezentace (beamer)
% Includuje se z ucitel/lekceXX/prezentace.tex jako:
%   % preamble-prez.tex -- sdílená preamble pro prezentace (beamer)
% Includuje se z ucitel/lekceXX/prezentace.tex jako:
%   % preamble-prez.tex -- sdílená preamble pro prezentace (beamer)
% Includuje se z ucitel/lekceXX/prezentace.tex jako:
%   \input{../spolecne/preamble-prez.tex}

\documentclass[aspectratio=169]{beamer}

% --- Jazyk a font ---
\usepackage{polyglossia}
\setmainlanguage{czech}
\usepackage{fontspec}

% --- Téma ---
\usetheme{default}
\usecolortheme{default}

\definecolor{microbit-blue}{HTML}{1E4D8C}
\definecolor{microbit-lightblue}{HTML}{D6E4F7}
\definecolor{code-bg}{HTML}{F0F0F0}

\setbeamercolor{frametitle}{bg=microbit-blue, fg=white}
\setbeamercolor{title}{fg=microbit-blue}
\setbeamercolor{block title}{bg=microbit-blue, fg=white}
\setbeamercolor{block body}{bg=microbit-lightblue}
\setbeamercolor{itemize item}{fg=microbit-blue}
\setbeamercolor{itemize subitem}{fg=microbit-blue!70}
\setbeamerfont{frametitle}{size=\large, series=\bfseries}
\setbeamertemplate{navigation symbols}{}
\setbeamertemplate{footline}{%
  \leavevmode%
  \hbox{%
    \begin{beamercolorbox}[wd=.5\paperwidth,ht=2.5ex,dp=1.125ex,leftskip=6pt]{author in head/foot}%
      \usebeamerfont{author in head/foot}\textcolor{gray}{\small Úvod do Pythonu na Micro:bitu}
    \end{beamercolorbox}%
    \begin{beamercolorbox}[wd=.5\paperwidth,ht=2.5ex,dp=1.125ex,rightskip=6pt,right]{date in head/foot}%
      \usebeamerfont{date in head/foot}\textcolor{gray}{\small\insertframenumber\,/\,\inserttotalframenumber}
    \end{beamercolorbox}}%
}

% --- Česky korektní uvozovky ---
\usepackage{csquotes}
\newcommand{\uv}[1]{\enquote{#1}}

% --- Zdrojový kód (Python) ---
\usepackage{listings}
\lstset{
  language=Python,
  basicstyle=\ttfamily\small,
  backgroundcolor=\color{code-bg},
  frame=single,
  framerule=0pt,
  xleftmargin=6pt,
  xrightmargin=6pt,
  breaklines=true,
  showstringspaces=false,
  keywordstyle=\color{microbit-blue}\bfseries,
  stringstyle=\color{teal},
  commentstyle=\color{gray}\itshape,
  literate=
    {á}{{\'a}}1 {é}{{\'e}}1 {í}{{\'i}}1 {ó}{{\'o}}1 {ú}{{\'u}}1
    {ů}{{\r{u}}}1 {č}{{\v{c}}}1 {š}{{\v{s}}}1 {ž}{{\v{z}}}1
    {Á}{{\'A}}1 {É}{{\'E}}1 {Í}{{\'I}}1 {Ó}{{\'O}}1 {Ú}{{\'U}}1
    {Č}{{\v{C}}}1 {Š}{{\v{S}}}1 {Ž}{{\v{Z}}}1 {ň}{{\v{n}}}1 {ř}{{\v{r}}}1,
}

% --- Zvýraznění pojmů ---
\newcommand{\pojem}[1]{\textbf{\textcolor{microbit-blue}{#1}}}
\newcommand{\pyinline}[1]{\texttt{#1}}

% --- Grafika a diagramy ---
\usepackage{tikz}

% --- Ostatní ---
\usepackage{graphicx}
\usepackage{hyperref}
\hypersetup{colorlinks=true, urlcolor=microbit-blue}


\documentclass[aspectratio=169]{beamer}

% --- Jazyk a font ---
\usepackage{polyglossia}
\setmainlanguage{czech}
\usepackage{fontspec}

% --- Téma ---
\usetheme{default}
\usecolortheme{default}

\definecolor{microbit-blue}{HTML}{1E4D8C}
\definecolor{microbit-lightblue}{HTML}{D6E4F7}
\definecolor{code-bg}{HTML}{F0F0F0}

\setbeamercolor{frametitle}{bg=microbit-blue, fg=white}
\setbeamercolor{title}{fg=microbit-blue}
\setbeamercolor{block title}{bg=microbit-blue, fg=white}
\setbeamercolor{block body}{bg=microbit-lightblue}
\setbeamercolor{itemize item}{fg=microbit-blue}
\setbeamercolor{itemize subitem}{fg=microbit-blue!70}
\setbeamerfont{frametitle}{size=\large, series=\bfseries}
\setbeamertemplate{navigation symbols}{}
\setbeamertemplate{footline}{%
  \leavevmode%
  \hbox{%
    \begin{beamercolorbox}[wd=.5\paperwidth,ht=2.5ex,dp=1.125ex,leftskip=6pt]{author in head/foot}%
      \usebeamerfont{author in head/foot}\textcolor{gray}{\small Úvod do Pythonu na Micro:bitu}
    \end{beamercolorbox}%
    \begin{beamercolorbox}[wd=.5\paperwidth,ht=2.5ex,dp=1.125ex,rightskip=6pt,right]{date in head/foot}%
      \usebeamerfont{date in head/foot}\textcolor{gray}{\small\insertframenumber\,/\,\inserttotalframenumber}
    \end{beamercolorbox}}%
}

% --- Česky korektní uvozovky ---
\usepackage{csquotes}
\newcommand{\uv}[1]{\enquote{#1}}

% --- Zdrojový kód (Python) ---
\usepackage{listings}
\lstset{
  language=Python,
  basicstyle=\ttfamily\small,
  backgroundcolor=\color{code-bg},
  frame=single,
  framerule=0pt,
  xleftmargin=6pt,
  xrightmargin=6pt,
  breaklines=true,
  showstringspaces=false,
  keywordstyle=\color{microbit-blue}\bfseries,
  stringstyle=\color{teal},
  commentstyle=\color{gray}\itshape,
  literate=
    {á}{{\'a}}1 {é}{{\'e}}1 {í}{{\'i}}1 {ó}{{\'o}}1 {ú}{{\'u}}1
    {ů}{{\r{u}}}1 {č}{{\v{c}}}1 {š}{{\v{s}}}1 {ž}{{\v{z}}}1
    {Á}{{\'A}}1 {É}{{\'E}}1 {Í}{{\'I}}1 {Ó}{{\'O}}1 {Ú}{{\'U}}1
    {Č}{{\v{C}}}1 {Š}{{\v{S}}}1 {Ž}{{\v{Z}}}1 {ň}{{\v{n}}}1 {ř}{{\v{r}}}1,
}

% --- Zvýraznění pojmů ---
\newcommand{\pojem}[1]{\textbf{\textcolor{microbit-blue}{#1}}}
\newcommand{\pyinline}[1]{\texttt{#1}}

% --- Grafika a diagramy ---
\usepackage{tikz}

% --- Ostatní ---
\usepackage{graphicx}
\usepackage{hyperref}
\hypersetup{colorlinks=true, urlcolor=microbit-blue}


\documentclass[aspectratio=169]{beamer}

% --- Jazyk a font ---
\usepackage{polyglossia}
\setmainlanguage{czech}
\usepackage{fontspec}

% --- Téma ---
\usetheme{default}
\usecolortheme{default}

\definecolor{microbit-blue}{HTML}{1E4D8C}
\definecolor{microbit-lightblue}{HTML}{D6E4F7}
\definecolor{code-bg}{HTML}{F0F0F0}

\setbeamercolor{frametitle}{bg=microbit-blue, fg=white}
\setbeamercolor{title}{fg=microbit-blue}
\setbeamercolor{block title}{bg=microbit-blue, fg=white}
\setbeamercolor{block body}{bg=microbit-lightblue}
\setbeamercolor{itemize item}{fg=microbit-blue}
\setbeamercolor{itemize subitem}{fg=microbit-blue!70}
\setbeamerfont{frametitle}{size=\large, series=\bfseries}
\setbeamertemplate{navigation symbols}{}
\setbeamertemplate{footline}{%
  \leavevmode%
  \hbox{%
    \begin{beamercolorbox}[wd=.5\paperwidth,ht=2.5ex,dp=1.125ex,leftskip=6pt]{author in head/foot}%
      \usebeamerfont{author in head/foot}\textcolor{gray}{\small Úvod do Pythonu na Micro:bitu}
    \end{beamercolorbox}%
    \begin{beamercolorbox}[wd=.5\paperwidth,ht=2.5ex,dp=1.125ex,rightskip=6pt,right]{date in head/foot}%
      \usebeamerfont{date in head/foot}\textcolor{gray}{\small\insertframenumber\,/\,\inserttotalframenumber}
    \end{beamercolorbox}}%
}

% --- Česky korektní uvozovky ---
\usepackage{csquotes}
\newcommand{\uv}[1]{\enquote{#1}}

% --- Zdrojový kód (Python) ---
\usepackage{listings}
\lstset{
  language=Python,
  basicstyle=\ttfamily\small,
  backgroundcolor=\color{code-bg},
  frame=single,
  framerule=0pt,
  xleftmargin=6pt,
  xrightmargin=6pt,
  breaklines=true,
  showstringspaces=false,
  keywordstyle=\color{microbit-blue}\bfseries,
  stringstyle=\color{teal},
  commentstyle=\color{gray}\itshape,
  literate=
    {á}{{\'a}}1 {é}{{\'e}}1 {í}{{\'i}}1 {ó}{{\'o}}1 {ú}{{\'u}}1
    {ů}{{\r{u}}}1 {č}{{\v{c}}}1 {š}{{\v{s}}}1 {ž}{{\v{z}}}1
    {Á}{{\'A}}1 {É}{{\'E}}1 {Í}{{\'I}}1 {Ó}{{\'O}}1 {Ú}{{\'U}}1
    {Č}{{\v{C}}}1 {Š}{{\v{S}}}1 {Ž}{{\v{Z}}}1 {ň}{{\v{n}}}1 {ř}{{\v{r}}}1,
}

% --- Zvýraznění pojmů ---
\newcommand{\pojem}[1]{\textbf{\textcolor{microbit-blue}{#1}}}
\newcommand{\pyinline}[1]{\texttt{#1}}

% --- Grafika a diagramy ---
\usepackage{tikz}

% --- Ostatní ---
\usepackage{graphicx}
\usepackage{hyperref}
\hypersetup{colorlinks=true, urlcolor=microbit-blue}


\title{Lekce 02: Proměnné a typy}
\subtitle{Úvod do Pythonu na Micro:bitu}
\date{}

\begin{document}

\begin{frame}
  \titlepage
\end{frame}

% ------------------------------------------------------------------
\begin{frame}{Co je proměnná?}
  \begin{columns}[T]
    \begin{column}{0.52\textwidth}
      \pojem{Proměnná} = krabice s nálepkou.

      \bigskip
      \begin{itemize}
        \item \textbf{Nálepka} = jméno proměnné
        \item \textbf{Obsah} = hodnota
      \end{itemize}

      \bigskip
      Kdykoli použijeme jméno,\\
      Python sáhne do krabice\\
      a vezme obsah.

      \bigskip
      \textcolor{gray}{\small
        Hodnotu v krabici lze kdykoli\\
        vyměnit — proto \emph{pro\-měn\-ná}.
      }
    \end{column}
    \begin{column}{0.44\textwidth}
      \centering
      \begin{tikzpicture}
        % krabice
        \draw[thick, microbit-blue, fill=microbit-lightblue, rounded corners=4pt]
          (0,0) rectangle (3.4,1.8);
        % nálepka
        \draw[fill=white, draw=microbit-blue, rounded corners=2pt]
          (0.8,1.5) rectangle (2.6,2.1);
        \node[font=\small\ttfamily\bfseries, text=microbit-blue] at (1.7,1.8) {jmeno};
        % obsah
        \node[font=\normalsize\ttfamily, text=microbit-blue!80] at (1.7,0.9)
          {"Alice"};
        % šipka přiřazení
        \node[font=\small, text=gray] at (1.7,-0.5) {\texttt{jmeno = "Alice"}};
        \draw[->, gray, thick] (1.7,-0.25) -- (1.7,0.1);
      \end{tikzpicture}
    \end{column}
  \end{columns}
\end{frame}

% ------------------------------------------------------------------
\begin{frame}[fragile]{Přiřazení hodnoty}

  \begin{lstlisting}
jmeno = "Alice"
  \end{lstlisting}

  \bigskip
  Čteme \textbf{zprava doleva}:
  \begin{enumerate}
    \item Vezmi hodnotu \pyinline{"Alice"}
    \item Ulož ji do proměnné \pyinline{jmeno}
  \end{enumerate}

  \bigskip
  \begin{lstlisting}
from microbit import *

jmeno = "Alice"
display.scroll(jmeno)   # zobrazí: Alice
  \end{lstlisting}

  \pause\bigskip
  \begin{alertblock}{Výhoda}
    Chceš zobrazit jiné jméno? Změníš \textbf{jednu hodnotu} nahoře —
    zbytek programu zůstane stejný.
  \end{alertblock}
\end{frame}

% ------------------------------------------------------------------
\begin{frame}[fragile]{Typy hodnot}

  Každá hodnota v Pythonu má svůj \pojem{typ}.

  \bigskip
  \begin{columns}[T]
    \begin{column}{0.46\textwidth}
      \begin{block}{\texttt{str} — řetězec (text)}
        Vždy v \textbf{uvozovkách}.\\[4pt]
        \pyinline{"Alice"}\\
        \pyinline{"1.A"}\\
        \pyinline{"16"}  \textcolor{gray}{\small\quad ← tohle je text!}
      \end{block}
    \end{column}
    \begin{column}{0.46\textwidth}
      \begin{block}{\texttt{int} — celé číslo}
        Bez uvozovek.\\[4pt]
        \pyinline{16}\\
        \pyinline{2025}\\
        \pyinline{-5}
      \end{block}
    \end{column}
  \end{columns}

  \bigskip
  \begin{lstlisting}
jmeno = "Alice"   # str
vek   = 16        # int
  \end{lstlisting}

  \pause
  \begin{alertblock}{Pozor}
    \pyinline{"16"} $\neq$ \pyinline{16} \quad — jedno je text, druhé číslo.
  \end{alertblock}
\end{frame}

% ------------------------------------------------------------------
\begin{frame}[fragile]{Spojování řetězců}

  Řetězce lze spojovat operátorem \pyinline{+}:

  \begin{lstlisting}
jmeno = "Alice"
pozdrav = "Ahoj, " + jmeno + "!"
display.scroll(pozdrav)   # zobrazí: Ahoj, Alice!
  \end{lstlisting}

  \pause\bigskip
  Co když chceme přidat číslo?

  \begin{lstlisting}
vek = 16
display.scroll("Věk: " + vek)   # CHYBA!
  \end{lstlisting}

  \pause
  \begin{exampleblock}{Řešení: \texttt{str()}}
    \begin{lstlisting}
display.scroll("Věk: " + str(vek))   # OK
    \end{lstlisting}
    \pyinline{str()} převede číslo na text.
  \end{exampleblock}
\end{frame}

% ------------------------------------------------------------------
\begin{frame}[fragile]{Digitální vizitka}

  \begin{lstlisting}
from microbit import *

jmeno = "Alice"
trida = "1.A"
hobby = "fotbal"

zprava = jmeno + ", " + trida + ", " + hobby
display.scroll(zprava)
  \end{lstlisting}

  \bigskip
  \begin{itemize}
    \item Proměnné definujeme \textbf{na začátku} — přehledné, snadno měnitelné.
    \item Zprávu sestavíme \textbf{spojením} — program nezávisí na konkrétních hodnotách.
  \end{itemize}
\end{frame}

% ------------------------------------------------------------------
\begin{frame}{Co jsme se naučili}
  \begin{itemize}
    \item[\checkmark] Proměnná = pojmenované úložiště hodnoty
    \item[\checkmark] Přiřazení: \pyinline{jmeno = "Alice"} — čteme zprava doleva
    \item[\checkmark] Typ \pyinline{str} — text v uvozovkách
    \item[\checkmark] Typ \pyinline{int} — celé číslo bez uvozovek
    \item[\checkmark] \pyinline{"16"} $\neq$ \pyinline{16}
    \item[\checkmark] Spojování řetězců operátorem \pyinline{+}
    \item[\checkmark] Převod čísla na text: \pyinline{str(vek)}
  \end{itemize}
\end{frame}

\end{document}
